\chapter{Related Work}
\label{cha:RelatedWork}
This chapter provides a comprehensive review of prior work across three key dimensions relevant for the development of a scientific QA system capable of providing cited text spans as evidence for its generated answer. First, it examines existing scientific QA datasets with respect to their suitability for SQuAI. It then compares different established cited text span extraction paradigms and their tradeoffs in terms of precision, efficiency and computational requirements before discussing automatic approaches for evaluating text span attribution quality in long-form, multi-hop settings. 
\section{Datasets}
In order to improve upon the existing SQuAI system, a highly specialized dataset of open-domain, long form scientific questions is needed.  Due to the scarcity and high cost of curating such datasets, only a limited number of suitable options are available. With SQuAI being designed as a system to answer complex scientific inquiries based on the unarXive 2022 Dataset \cite{unarXive2022}, further requirements are introduced, namely: 
\begin{enumerate}
    \item Domain Alignment: Questions must fall under the knowledge scope of the underlying knowledgebase, in order for the system to be able to provide cited text spans as proof alongside its answer. In terms of the SQuAI System, this boils down to questions targeting physics, mathematics and computer science that are also part of the papers in the unarXive 2022 Dataset.
    \item High Complexity: Questions must go beyond simple factoid questions and require complex reasoning capabilities to answer, as questions posed by scientists usually require synthesizing, interpreting and evaluating data, requiring high cognitive abilities to interact with. In order to be able to test if the SQuAI system can be used under such circumstances, the questions in the dataset must also have a similar degree of complexity.
    \item Multi-Hop: Questions must require the information in multiple documents to be answered, in order to test the ability of the system to synthesize information from multiple sources, as is required in naturally occurring scientific research.
    \item Decontextualization: Questions must be answerable in an open-domain setting and cannot contain any kind of reference to specific papers and their contents that lose meaning or become ambiguous outside of the context of their paper. As a result, questions must contain enough context such that they are standalone answerable.
    \item Text format: Questions can only target information in a paper that is present in the form of text, as current versions of the SQuAI system do solely rely on textual information, not taking into account images, graphs or similar. 
    \item Ground Truth: In order to evaluate the correctness of the model's answer, a dataset must provide a ground answer alongside each question.
\end{enumerate}

However, comparing some of the most prominent datasets shows that none fits all discussed criteria, finally resulting in the need for a newly curated dataset.
% With each of the existing Datasets failing to meet those reqirements, this leaves us with no choice but to curate our own.
\begin{table}[H]
\centering
\label{tab
:dataset_comparison}
\resizebox{\textwidth}{!}{%
% \begin{tabular}{lccccccc}
%     \hline
%     \textbf{Dataset} & 
%     \textbf{Domain Alignment} & \textbf{High Complexity} & \textbf{Multi-Hop} & 
%     \textbf{Decontextualization} & \textbf{Text format}  & \textbf{Ground Truth} & \textbf{Dataset size} \\ \hline
%     Qasper      \cite{qasper}           & neutral (NLP only)           & shallow      & no              & no        & yes             & yes                               & \\
%     ALCE        \cite{alceBenchmark}    & no (General Knowledge)       & shallow      & yes             & yes       & yes             & yes (unreliable)                  & \\
%     SPIQA       \cite{spiqa}            & yes (CS)                     & complex      & no              & partially & no (multimodal) & yes (visual representation)        &\\             
%     QASA        \cite{qasa}             & yes (CS)                     & complex      & no              & no        & yes             & yes (multi-layered)               &\\ 
%     ExpertQA    \cite{expertQA}         & yes (broad spectrum)         & very complex & yes             & yes       & yes             & yes (but evidence are weblinks)   &  too few if filtered to squai topics  \\                                 
%     SciDQA      \cite{scidqa}           & yes (ML)                     & very complex & no (only 11\%)  & no        & no              & yes (with evidence)               & \\                                                             
%     PeerQA      \cite{peerqa}           & yes (ML \& NLP)              & very complex & no              & no        & yes             & yes (with evidence)               & \\   \hline                                                    
% \end{tabular}%
\begin{tabular}{lccccccc}
    \hline
    \textbf{Dataset} & 
    \textbf{\makecell{Domain\\Alignment}} & \textbf{\makecell{High\\Complexity}} & \textbf{Multi-Hop} & 
    \textbf{\makecell{Decontextual-\\ization}} & \textbf{\makecell{Text\\format}}  & \textbf{\makecell{Ground\\Truth}} & \textbf{\makecell{Sufficient Dataset\\size (>500)}} \\ \hline
    Qasper      \cite{qasper}           & partial           & shallow      & no              & no        & yes             & yes    & yes \\
    ALCE        \cite{alceBenchmark}    & no                & shallow      & yes             & yes       & yes             & yes    & yes \\
    SPIQA       \cite{spiqa}            & yes               & complex      & no              & partially & no & yes    & yes \\             
    QASA        \cite{qasa}             & yes               & complex      & no              & no        & yes             & yes    & yes \\ 
    ExpertQA    \cite{expertQA}         & yes               & very complex & yes             & yes       & yes             & yes    & no \tablefootnote{question-answer pairs are broadly distributed across diverse fields of study, resulting in insufficient topical overlap with unarXive 2022}   \\                                 
    SciDQA      \cite{scidqa}           & yes               & very complex & no              & no        & no              & yes    & yes \\                                                             
    PeerQA      \cite{peerqa}           & yes               & very complex & no              & no        & yes             & yes.   & yes \\   \hline                                                    
\end{tabular}%
}
% TODO better formattig here
    \caption{Comparison of Scientific QA Datasets against SQuAI Requirements\\ 
    Abbreviations: CS = Computer Science; ML = Machine Learning; NLP = Natural Language Processing.}
\end{table}

\section{Cited Text Span Extraction}
Context extraction methods can be most broadly divided into pre- and post-hoc extraction methods, depending on the time the context attribution happens. While pre-hoc systems such as QASPER \cite{qasper} and QASA \cite{qasa} extract citable passages before and during runtime of the question-answering system, post-hoc systems search and attach supporting evidence afterwards. Post-hoc systems have been shown to outperform pre-hoc systems, achieving higher citation coverage and answer correctness as well as lowering model hallucination \cite{generationTimeVsPostHocCitation}. 
While a number of systems rely on  zero-shot LLM solutions for context extraction \cite{onSynthesizingDataForContextAttributionInQuestionAnswering} \cite{whereDidThatComeFrom} \cite{ContextBasedQuotationRecommendation}, work on the DECOMPTUNE \cite{decompositionEnhancedTraining} system has shown that decomposing facts and decontextualizing answer sentences before attributing source text spans with LLM prompting can yield significant improvements, especially when applied to answers for long-form, multi-hop questions. In order to eliminate the probability of hallucinated model output, models can be finetuned \cite{citationContextExtractionUsingTransformerBasedModel} or the answer space restricted to the start and end token of the cited text span, forcing strict citation \cite{citationContextExtractionUsingATransformerBasedModel}.
% check if they are really all zero-shot LLM solutions
% shit sentence 

% Instead of treating the problem as a retrieval task the DECOMPTUNE \cite{decompositionEnhancedTraining} system approaches post-hoc attribution from a reasoning perspective. The authors observed that decomposing complex answers into logical units before attributing context yields significant improvements in citatioxn quality compared to zero -hot LLM solutions.

So-called white-box methods rely on access to the model's internal weights, rendering them incompatible with proprietary LLMs. The LAQuer \cite{laquer} system, as well as earlier token-embedding approaches it is built upon \cite{peeringIntoTheMindOfLanguageModels}, analyzes internal hidden states to extract the most likely text span that a user-highlighted part of the LLM's answer is based upon. To pinpoint a certain text span, the system calculates the maximum cosine similarity between individual tokens in the source text and the averaged hidden representation of the highlighted span, for which attribution is sought. The system then extracts the precise quote from the surrounding window of text. However, this method has been shown by the authors to perform significantly worse than LLM Prompting, with drastically lower accuracy and missing suitable citations, especially in larger search spaces, e.g. longer text. 
Similarly, AttnAttribute \cite{onMechanisticCircuitsForExtractiveQuestions} identifies "Attention Heads", tracking the model's internal attention values during generation time. Just like LAQuer, it then extracts text spans from the passage surrounding the anchor with the highest model attention. 
Benchmarks made as part of the development of DECOMPTUNE \cite{decompositionEnhancedTraining} concluded that both of those methods performed significantly worse than LLM Prompting.

% double check qasper and case claim
%verify "zero-shot LLM based context extraction"  really true for all 3
% highlight how costly some of them are
Recent advancements have explored generative approaches. Rather than retrieving and reranking text spans, such systems generate them. In order to prevent hallucinations, the output vocabulary of the systems is typically restricted to the underlying retrieval corpus by applying constrained decoding mechanisms, such as prefix trees or FM-indices \cite{autoregressiveSeachEnginesSeal}. 
GERE \cite{generativeEvidenceRetrievalForFactVerification}, as one such system using a prefix tree, applies this concept of restricted output in the document retrieval step, generating document titles for selecting appropriate source texts, as well as sentence identifiers mapped to individual sentences. This system, while removing the need for a massive document index as well as showing significant improvements in precise and correct text span extraction compared to baselines such as Dense Passage Retrieval (DPR) and BM25, requires the prefix tree to be precomputed and stored in memory, making it highly inflexible for changes in the underlying corpus.  
Similarly, MRR-FV \cite{mrr-fv} further improves this approach with multi-hop reasoning capabilities. Instead of generating sentence identifiers, it generates sentences directly, leveraging the natural language understanding capabilities of LLMs, leading to improvements in precision and quality of attributed spans, but also further bloating the prefix tree, making the system even harder to scale.
% ell overall improvements in span preciness and quality.  yeah this should be better explained

Despite the diversity of proposed paradigms, current work mostly evaluates them in isolation, restricting comparison to baselines or direct predecessors on heterogeneous tasks, making broad comparisons for practical applicability unfeasible. Consequently, this work introduces a hierarchy of extraction methods, increasing in both reasoning capabilities and computational requirements, benchmarking both accuracy and efficiency on a unified framework, allowing for systematic trade-off analysis between paradigms in the context of open-domain multi-hop scientific QA.
% restricting wrong word
% wrap this up with what my thesis adds, what is the contribution and why
    % practical limitations- scalability, run time constraints, reliance on interal model weights, heacy prefix treees
    % spectrum of retrieval methods - unified benchmark needed
    % compared in isolation, task specific
    % methods in isolation, not in composition

\section{Span Evaluation}
While human annotation has always been the gold standard for assessing natural language generation (NLG) tasks, relying solely on it is a bottleneck for developing models of increasing size and complexity. While providing feedback with extensive reasoning and natural language understanding, human evaluation suffers from being costly, time-intensive and, depending on the specific task, highly subjective. Thus, lowering or even replacing the need for human evaluation completely is an important problem. 

Early metrics heavily rely on measuring the lexical overlap between text spans, for instance of a human-written ground truth and an LLM-generated answer for the same questions. The most prominent examples include BLEU \cite{bleu} and ROUGE \cite{rouge}, both being n-gram based metrics. While the former is precision-oriented, measuring how many of its n-grams also appear in the human reference span, the latter applies a recall-oriented approach, measuring how many n-grams of the human reference are also present in the LLM-generated text span. Being computationally inexpensive and language independent, these metrics have established themselves as rigid industry standards and baselines for other methods. Their alignment with human judgement, however, has been shown to be low, suffering from the lack of semantic understanding, leading to penalizing valid paraphrasing as well as missing contradictory or hallucinated content.

Partly overcoming those limitations, embedding-based metrics like BERTScore \cite{bertScore} and MoverScore \cite{moverScore} emerged, measuring semantic instead of lexical similarity. BERTScore does one-to-one greedy matching on vector embeddings, providing precision, recall and F1 score metrics, just like BLEU and ROUGE, but on an embedding level. MoverScore improves upon this idea, utilizing Earth Mover's Distance (EMD) to allow for introducing one-to-many relationships, allowing for matching of single words with phrases, so-called soft matching, as well as introducing Inverse Document Frequency (IDF) for enhanced lexical matching, to better reflect the importance of specific keywords to be included. Though significantly more expensive to compute, those metrics show significant improvements in semantic robustness as well as intrinsic context awareness, due to their vector embeddings. However, both of those metrics still require a reference answer to compare to. BartScore \cite{bartScore} eliminates this requirement, measuring the conditional probability that one text span is generated from another using an encoder-decoder model, arguing that the probability is higher if the claims in the answer are also deeply rooted in the source text.

Besides similarity-based metrics, other research suggests approaching the task from a logical entailment perspective. The AIS \cite{measuringAttributionInNaturalLanguageGenerationModels} framework proposed a unified conceptual definition for such, arguing that a statement is considered Attributable to Identified Sources (AIS) iff completely corroborated by a source. In order to verify this, the authors imposed additional constraints being 1) that the claim made in an answer sentence must be standalone interpretable, thus decontextualized and 2) the information in the text span must support all the information present in the answer span it is attributed to. Previously discussed similarity-based metrics like BLEU, ROUGE or token-level F1, as well as solely relying on an out-of-the-box Natural Language Inference (NLI) model, have been shown to have low correlation with human judgement for both factual consistency and logical entailment assessment \cite{TRUE} \cite{evaluatingFactualConsistencyOfAbstractiveTextSummarization}. 
To address these limitations, systems relying on NLI models like FactCC \cite{evaluatingFactualConsistencyOfAbstractiveTextSummarization} argued that evaluating logical entailment on sentence-to-sentence level is insufficient, since information in sources might be spread across multiple passages. As a result, these systems increased their evaluation scope to the whole source text, leading to better judgement capabilities, but also increase in model input. The authors of FactCC observed that a majority of misclassifications were commonsense mistakes; the TRUE benchmark \cite{TRUE} observed improvements with increasing model size but also degradation with longer input, which is further supported by research on the SummaC \cite{summaC} framework.
To shrink input size, SummaC increased evaluation granularity by dividing both the source text and the model answer into smaller units, checking each combination of them for entailment and combining the obtained scores. 

The Q\textsuperscript{2} system  \cite{q2evaluatingFactualConsistencyInKnowledgeGroundedDialogues} proposed using a Question-Generation - Question-Answering (QG-QA) method, generating questions to which the LLM-generated answer is the solution. Those questions are then answered with the original source document and lastly the so obtained answers are then compared with the original LLM answer using an NLI model. The authors presented significant improvements compared to baseline methods, further enhanced by the TRUE benchmark by applying an ensemble approach utilizing ANLI (standard NLI), SummaC (increased granularity NLI) as well as the Q\textsuperscript{2} (QG-QA) method. The ensemble was shown to beat every standalone evaluation metric previously discussed on all datasets benched in TRUE, emphasizing the need for combined, complementary evaluation metrics.

Recently, the evaluation paradigm has shifted to LLM-as-a-judge approaches. While early versions like GPTScore \cite{gptScore} tried evaluating based on model internals, using the conditional token probability, thus the confidence of a model in a given answer, they have quickly been shown to be much less effective than prompting the LLM directly to grade \cite{gval}. The performance of LLMs in evaluation tasks can additionally be enhanced by applying a Chain-of-thought approach, prompting the model to break down complex tasks into simpler ones and adding intermediate reasoning steps. Closely resembling human cognitive processes, this method yields significant improvements, especially with larger models and problems of increasing complexity \cite{chainOfThoughtPrompting}, reaching agreement rates comparable to human judges for open-ended multi-turn dialogues \cite{judgingLLMAsAJudgeMTBench}, thus eliminating human annotation constraints. 

When applied to long-form text containing multiple pieces of information, binary judgements become insufficient as they cannot reflect partial attribution. FactScore \cite{factScore} addresses this by breaking down answers into atomic facts using an LLM, obtaining a percentage measure of supportedness. AttrScore \cite{automaticEvaluationOfAttribution} introduces additional evaluation criteria, such as measures for attribution, extrapolation and contradiction, in order to further characterize the type of attribution error. Similarly, CiteEval \cite{citeEvalPrincipleDrivenCitationEvaluationForSourceAttribution} introduces penalties for attributions containing unnecessary noise and redundant citations while advocating for evaluation against full source documents, in order to not lose information in the surrounding context. The RAGAS \cite{ragas} framework operationalized these findings, providing a suite of automated metrics like faithfulness, answer relevance and context relevance, without requiring human annotation. However, the importance of penalizing undesired behaviors is further emphasized by the MT-Benchmark \cite{judgingLLMAsAJudgeMTBench}, showing that LLMs have a multitude of biases, such as favoring longer and LLM-generated text spans or switching preferences between answers when presented in a different order. Furthermore, LLMs perform remarkably poorly with fine-grained information, with a majority of mistakes being caused by overlooking numbers and entities or hallucinating logical connections \cite{attributionBench}. Using increasingly large models additionally introduces high computational cost requirements. To address this, systems like ARES \cite{ares} rely on LLM-generated, domain-specific, synthetic training data to finetune smaller models. Using Prediction-Powered-Inference (PPI), ARES leverages a small human-annotated dataset to calculate the LLM's annotation error on the synthetic data, the "rectifier", to provide statistical confidence bounds and enhance its evaluation, beating existing zero-shot solutions like RAGAS.  

Automatic evaluation of cited text span attribution remains a challenging task. Despite the wide range of approaches, no single metric being capable of reliably and comprehensively assessing attribution performance as a whole has emerged. Increasingly complex tasks, such as multi-hop questions and challenging application domains like scientific QA, introduce further evaluation dimensions. Consequently, combining multiple evaluation metrics with complementary strengths becomes necessary to measure task-specific aspects within a unified framework.

% wide range of methods
% remains challenging task
% no single metrics is best
% not only binary judgemnts

% wtf supportedness
% meh -  in order to not lose information in the surrounding context. 
% todo talk about ragas
% 
% 


% starting with the LLM part here...
% CiteEval \cite{citeEvalPrincipleDrivenCitationEvaluationForSourceAttribution} 


% --> Tobias: verbindung zu eigener arbeit als letzter satz


% human annotiation costly, time-intensive, highly subjective
% as intrinsic context awareness yeah think about that again, should cite some proof for my claims